\documentclass{article}
\usepackage{graphicx} % Required for inserting images
\usepackage{amsfonts}
\usepackage{amsmath}

\title{texture synthesis (Gaussian splatting part)}
\author{Jianfeng Pang }
\date{April 2026}

\begin{document}

\maketitle

\section{Introduction}

\section{Related work}


\section{Method}



\subsection{Mesh-guided Gaussian Splatting from Synthesized Textures}
The preceding reference-guided texture synthesis stage produces a textured mesh by projecting the refined multi-view images back to the UV domain. We denote the resulting asset as
\begin{equation}
\hat{\mathcal{M}} = (\mathcal{V}, \mathcal{F}, \mathcal{T}),
\end{equation}
where $\mathcal{V}=\{v_j\}_{j=1}^{|\mathcal{V}|}$ is the set of vertices, $\mathcal{F}$ is the set of triangular faces, and $\mathcal{T}$ is the synthesized UV texture map. From $\hat{\mathcal{M}}$, we learn a mesh-guided Gaussian representation $\mathcal{G}$ for novel-view rendering and editable rendering. Instead of treating Gaussian Splatting as an unconstrained point-based representation \cite{kerbl3Dgaussians}, we bind the Gaussian splats to the mesh faces following the mesh-based parameterization of GaMeS \cite{waczynska2024games}, so that the learned radiance field remains aligned with the synthesized surface texture.

\subsubsection{Multi-view supervision construction}
Let $\mathcal{C}=\{c_i\}_{i=1}^{N}$ be the set of cameras used for Gaussian optimization. Each camera is represented as
\begin{equation}
c_i=(K_i,E_i),
\end{equation}
where $K_i$ denotes the intrinsic matrix and $E_i=[R_i|t_i]$ denotes the world-to-camera extrinsic transform. The camera poses and intrinsics are read from the NeRF-style files \texttt{transforms\_train.json} and \texttt{transforms\_test.json} \cite{mildenhall2020nerf}. Each frame stores a camera-to-world pose and an image path, while the horizontal field of view is converted to the focal length used during Gaussian rasterization.

For each supervision camera $c_i$, we construct a target image $I_i \in \mathbb{R}^{H\times W\times 3}$. When the final UV texture is available, the target view is rendered from the textured mesh:
\begin{equation}
I_i=\mathbb{R}_{\mathrm{mesh}}(\hat{\mathcal{M}},c_i).
\end{equation}
Equivalently, if the previous stage stores view-wise synthesized images, $I_i$ can be taken directly from the corresponding refined or inpainted output. This gives the Gaussian training set
\begin{equation}
\mathcal{D}_{\mathrm{GS}}=\{(I_i,c_i)\}_{i=1}^{N}.
\end{equation}
The training and test cameras are kept separate. The cameras in \texttt{transforms\_train.json} define the photometric supervision, while \texttt{transforms\_test.json} is used for held-out or modified camera-view rendering. Therefore, changing the test poses evaluates the learned representation under new camera coordinates rather than forcing it to reproduce the original synthesized images.

\subsubsection{Mesh-bound Gaussian parameterization}
We define the Gaussian representation as
\begin{equation}
\mathcal{G}=\{g_{f,k}\mid f\in\mathcal{F},\;k=1,\ldots,K\},
\end{equation}
where $K$ Gaussian splats are initialized on every mesh face. Following GaMeS \cite{waczynska2024games}, for a triangular face $f=(v_{f,1},v_{f,2},v_{f,3})$, the center of the $k$-th Gaussian is written in barycentric coordinates:
\begin{equation}
x_{f,k}
=
\alpha_{f,k,1}v_{f,1}
+\alpha_{f,k,2}v_{f,2}
+\alpha_{f,k,3}v_{f,3},
\end{equation}
with the simplex constraint
\begin{equation}
\alpha_{f,k,r}\ge 0,\qquad
\sum_{r=1}^{3}\alpha_{f,k,r}=1.
\end{equation}
In practice, the barycentric weights are obtained by normalizing unconstrained learnable parameters. This constraint keeps the Gaussian centers on the mesh surface and prevents the representation from drifting into unrelated free space.

As in 3D Gaussian Splatting \cite{kerbl3Dgaussians}, each Gaussian contains a center, opacity, anisotropic scale, rotation, and color coefficients:
\begin{equation}
g_{f,k}=
\left(
x_{f,k},
\Sigma_{f,k},
o_{f,k},
\mathbf{h}_{f,k}
\right),
\end{equation}
where $o_{f,k}\in[0,1]$ is opacity and $\mathbf{h}_{f,k}$ denotes RGB color or spherical harmonics coefficients. For view-dependent color, the emitted radiance along direction $d$ is
\begin{equation}
\mathbf{c}_{f,k}(d)
=
\sum_{\ell=0}^{L}\sum_{m=-\ell}^{\ell}
\mathbf{h}_{f,k}^{\ell m}Y_{\ell m}(d),
\end{equation}
where $Y_{\ell m}$ are spherical harmonics basis functions.

The covariance is also tied to the face geometry, following the mesh-face construction in GaMeS \cite{waczynska2024games}. Let the face normal be
\begin{equation}
n_f=
\frac{(v_{f,2}-v_{f,1})\times(v_{f,3}-v_{f,1})}
{\|(v_{f,2}-v_{f,1})\times(v_{f,3}-v_{f,1})\|_2}.
\end{equation}
We construct two tangent directions $t_{f,1}$ and $t_{f,2}$ by orthonormalizing the face edges, and form a local frame
\begin{equation}
R_f=[n_f,\;t_{f,1},\;t_{f,2}].
\end{equation}
The covariance is parameterized as
\begin{equation}
\Sigma_{f,k}
=
R_fS_{f,k}S_{f,k}^{T}R_f^{T},
\end{equation}
where
\begin{equation}
S_{f,k}
=
\mathrm{diag}(\epsilon,\rho_{f,k}s_{f,1},\rho_{f,k}s_{f,2}).
\end{equation}
Here $\epsilon$ is a small fixed scale in the normal direction, $s_{f,1}$ and $s_{f,2}$ are face-dependent tangent scales, and $\rho_{f,k}>0$ is a learnable multiplier. Thus the Gaussian is flat along the surface normal and anisotropic along the face tangents. Since both $x_{f,k}$ and $\Sigma_{f,k}$ are functions of the mesh vertices, editing the mesh automatically updates the Gaussian centers, orientations, and scales.

\subsubsection{Differentiable Gaussian rendering}
Given a camera $c_i=(K_i,E_i)$, all Gaussians are projected to the image plane using differentiable Gaussian rasterization \cite{kerbl3Dgaussians}. For the $j$-th Gaussian, where $j$ indexes a pair $(f,k)$, the projected mean is
\begin{equation}
\mu_{i,j}
=
\pi\!\left(K_iE_i x_j\right),
\end{equation}
where $\pi(\cdot)$ denotes perspective division. The 3D covariance is mapped to screen space using the local projection Jacobian $J_{i,j}$:
\begin{equation}
\Sigma'_{i,j}
=
J_{i,j}R_i\Sigma_jR_i^{T}J_{i,j}^{T}.
\end{equation}
For a pixel $p$, the opacity-weighted contribution of the projected Gaussian is
\begin{equation}
a_{i,j}(p)
=
o_j
\exp\!\left(
-\frac{1}{2}(p-\mu_{i,j})^T
(\Sigma'_{i,j})^{-1}
(p-\mu_{i,j})
\right).
\end{equation}
After depth ordering, differentiable alpha compositing produces the rendered image:
\begin{equation}
\hat{I}_i(p)
=
\sum_{j\in\mathcal{N}_i(p)}
T_{i,j}(p)\,
a_{i,j}(p)\,
\mathbf{c}_j(d_{i,j}),
\end{equation}
where $\mathcal{N}_i(p)$ is the set of splats overlapping pixel $p$, $d_{i,j}$ is the view direction, and
\begin{equation}
T_{i,j}(p)
=
\prod_{\ell<j}
\left(1-a_{i,\ell}(p)\right)
\end{equation}
is the accumulated transmittance before the $j$-th Gaussian. The renderer is differentiable with respect to opacity, spherical harmonics color, tangent scale, barycentric coordinates, and optionally mesh vertex positions.

\subsubsection{Optimization objective}
The mesh-guided Gaussian representation is optimized by comparing the rendered image $\hat{I}_i$ with the target image $I_i$. For a mini-batch of training cameras $\mathcal{B}$, we use
\begin{equation}
\mathcal{L}_{\mathrm{photo}}
=
\frac{1}{|\mathcal{B}|}
\sum_{i\in\mathcal{B}}
\left[
(1-\lambda)\mathcal{L}_{1}(\hat{I}_i,I_i)
+
\lambda\left(1-\mathrm{SSIM}(\hat{I}_i,I_i)\right)
\right],
\end{equation}
where $\lambda$ controls the relative weight of structural similarity. The optimized parameters are
\begin{equation}
\Theta=
\{\alpha_{f,k},o_{f,k},\rho_{f,k},\mathbf{h}_{f,k}\}_{f,k},
\end{equation}
and the vertex set $\mathcal{V}$ may also be optimized when small geometry corrections are desired. In our setting, the textured mesh provides a reliable geometric prior, so fixing the topology makes training stable and keeps the Gaussian representation aligned with the synthesized texture.

This differs from vanilla 3D Gaussian Splatting \cite{kerbl3Dgaussians}, where
\begin{equation}
\mathcal{G}_{\mathrm{free}}
=
\{(\mu_j,\Sigma_j,o_j,\mathbf{h}_j)\}_{j=1}^{M},
\qquad
\mu_j\in\mathbb{R}^{3},
\end{equation}
and the centers are unconstrained points in space. By constraining the Gaussian centers and covariances with mesh topology, our representation better preserves the geometry of the synthesized mesh and is more suitable for downstream mesh-level editing.

\subsubsection{Novel-view rendering and evaluation}
After training, a novel view is rendered by providing a query camera
\begin{equation}
c^{*}=(K^{*},E^{*}),
\end{equation}
and applying the same Gaussian rasterizer:
\begin{equation}
\hat{I}^{*}
=
\mathbb{R}_{\mathrm{GS}}(\mathcal{G},c^{*}).
\end{equation}
The output depends only on the learned Gaussian representation and the supplied camera parameters. Therefore, if the test camera file is modified, the rendered result follows the new camera coordinates and does not depend on the old ground-truth images.

When the test cameras match the target images, full-reference metrics such as PSNR, SSIM, and LPIPS can be computed between $\hat{I}_i$ and $I_i$. If the test cameras are changed to novel poses, these metrics no longer have strict correspondence because the target and rendered images describe different rays. We therefore also use geometry-aware criteria. Let $S_i^{\mathrm{mesh}}$ be the mesh silhouette rendered at $c_i$, and let $S_i^{\mathrm{GS}}$ be the foreground mask induced by accumulated Gaussian opacity. We define
\begin{equation}
\mathrm{Coverage}
=
\frac{|S_i^{\mathrm{GS}}\cap S_i^{\mathrm{mesh}}|}
{|S_i^{\mathrm{mesh}}|+\epsilon},
\end{equation}
\begin{equation}
\mathrm{Hole}
=
\frac{|S_i^{\mathrm{mesh}}\setminus S_i^{\mathrm{GS}}|}
{|S_i^{\mathrm{mesh}}|+\epsilon},
\qquad
\mathrm{Leakage}
=
\frac{|S_i^{\mathrm{GS}}\setminus S_i^{\mathrm{mesh}}|}
{|S_i^{\mathrm{GS}}|+\epsilon}.
\end{equation}
Foreground coverage, hole ratio, leakage ratio, and sharpness better indicate whether the mesh-guided Gaussian representation remains complete and compact under novel camera viewpoints.


\section{Experiments}

\section{Conclusion}

\bibliographystyle{plain}
\bibliography{reference}


\end{document}
